\documentclass[11pt]{article}

% =============================================================================
% layer_to_layer_architecture.tex
%
% Layer-to-layer architecture of the QED exact-diagonalization library, from the
% top-level Python API down to the rudimentary matvec kernels and backends.
% Reflects the post-"symmetry-on-the-engine" state (branch
% lanczos-unification-gen1-retire): one matvec engine, three orthogonal axes.
% =============================================================================

\usepackage[margin=1in]{geometry}
\usepackage{amsmath,amssymb}
\usepackage{listings}
\usepackage{xcolor}
\usepackage{booktabs}
\usepackage{enumitem}
\usepackage{hyperref}
\hypersetup{colorlinks=true,linkcolor=blue!50!black,urlcolor=blue!50!black}

\definecolor{kw}{rgb}{0.13,0.29,0.53}
\definecolor{cm}{rgb}{0.30,0.50,0.30}
\definecolor{ty}{rgb}{0.55,0.20,0.50}
\lstset{
  basicstyle=\ttfamily\footnotesize,
  keywordstyle=\color{kw}\bfseries,
  commentstyle=\color{cm}\itshape,
  stringstyle=\color{ty},
  showstringspaces=false,
  breaklines=true,
  columns=fullflexible,
  language=C++,
  morekeywords={constexpr,override,noexcept,nullptr,uint64_t,std,template,typename,struct,class,using},
  frame=single,
  framesep=4pt,
  rulecolor=\color{black!25},
  backgroundcolor=\color{black!2},
}

\newcommand{\code}[1]{\texttt{\small #1}}
\newcommand{\file}[1]{\texttt{\footnotesize #1}}

\title{\textbf{QED Exact Diagonalization:\\ Layer-to-Layer Architecture}\\[4pt]
\large From the Python API down to the matvec kernels}
\author{Architecture reference}
\date{\today}

\begin{document}
\maketitle

\begin{abstract}
This document describes the QED exact-diagonalization (ED) library
layer by layer, from the top-level Python entry points down to the
rudimentary sparse matrix--vector (SpMV) kernels and the linear-algebra
backends. The central organizing principle is that an ED computation is
one point in a \emph{three-axis cube}---\textbf{reduction} (which basis
the Hamiltonian acts on) $\times$ \textbf{method} (which eigensolver)
$\times$ \textbf{deployment} (which hardware backend)---and that these
three axes compose independently above a single seam, the
\code{MatVecOperator}. We trace each axis to its concrete code, show the
contracts between adjacent layers, and present the rudimentary
implementations: the term SoA, the gather/scatter SpMV kernels, the
basis policies (including the multi-target generalization that lets
non-abelian symmetry ride the same engine as everything else), and the
symmetry reduction (group closure, numerical irrep decomposition, and
the symmetry-adapted basis).
\end{abstract}

\tableofcontents

% =============================================================================
\section{The organizing model: three orthogonal axes}
% =============================================================================

A solve is fully specified by a point in
\[
\underbrace{\text{Full} \mid \text{Sz} \mid \text{abelian} \mid \text{non-abelian} \mid \text{Sz}\!\times\!\text{sym}}_{\textbf{A: reduction (basis)}}
\;\times\;
\underbrace{\text{dense} \mid \text{Lanczos} \mid \text{Krylov--Schur} \mid \text{FTLM/TPQ} \mid \text{cf-DSSF}}_{\textbf{B: method}}
\;\times\;
\underbrace{\text{CPU} \mid \text{GPU} \mid \text{MPI}}_{\textbf{C: deployment}}.
\]
The library realizes this cube so that the axes are independent: a
reduction is a \emph{basis policy}, a method is a \emph{kernel
consuming a matvec}, and a deployment is a \emph{backend template
parameter}. The one seam they meet at is the abstract operator
interface \code{ed::matvec::MatVecOperator}, which exposes only
$v \mapsto Hv$ and the dimension. Everything above that seam is
symmetry-blind; everything below it is method-blind.

% =============================================================================
\section{Physical layering: the static libraries}
% =============================================================================

The build is a small DAG of static libraries
(\file{cmake/EDLibraries.cmake}). Arrows mean ``links / depends on'';
lower is more fundamental.

\begin{center}
\begin{lstlisting}[language=,basicstyle=\ttfamily\scriptsize]
 _core   (pybind11 module: python/qed/_bindings/qed_bindings.cpp) -- links everything
   |
   +-- ed_cli        ed_solvers_gpu        ed_distributed     ed_input
   |    (verbs)      (CUDA: facade lanczos  (MPI/NCCL)        (HamiltonianBuilder)
   |                  kernel<CudaBackend>,
   |                  gpu_ftlm, kpm,
   |                  symmetry_adapted_gpu)
   |                       | PUBLIC
   v                       v
   +------------>  ed_solvers_cpu  (lanczos_kernel consumers, planner, thermal,
   |                |   |   |        ftlm, symmetry_adapted_solve)
   |                |   |   +--> ed_symmetry  (group, irreps, SAB/sector builders)
   |                |   +------> ed_matvec    (backends + basis policies + kernels)
   v                v
 ed_dssf          ed_core  (Operator, SubspaceOperator, FixedSzOperator, term SoA)
                    |
                  ed_io --> ed_parallel  (OpenMP / pthreads)
\end{lstlisting}
\end{center}

Two edges are load-bearing:
\begin{itemize}[nosep]
  \item \code{ed\_solvers\_cpu} $\to$ \code{ed\_symmetry} is one-way. Hence
  the symmetry reduction (which builds the basis) lives in
  \code{ed\_symmetry}, but anything that \emph{solves} (calls the Krylov
  kernel) lives one layer up in \code{ed\_solvers\_cpu}.
  \item \code{ed\_matvec} and \code{ed\_symmetry} are siblings (both
  $\to$ \code{ed\_core}). The basis \emph{policies} live in
  \code{ed\_matvec}; the group/irrep/SAB \emph{construction} lives in
  \code{ed\_symmetry}.
\end{itemize}

% =============================================================================
\section{Logical layering and the inter-layer contracts}
% =============================================================================

\begin{center}
\begin{lstlisting}[language=,basicstyle=\ttfamily\scriptsize]
 L0  Frontend     python/qed/*.py            solve / thermal / dssf / symmetry_adapted_*
 L1  Bindings     qed_bindings.cpp           Python<->C++; use_gpu, n_up, method
 ----------------------------------------------------------------------------------
 L2  Consumers    solvers/, thermal/, dssf/  loop sectors -> pick method -> recombine x d_G
 L3  Method       krylov/, lanczos.cpp       dense / Lanczos / FTLM / cf-DSSF
 ===== SEAM:  MatVecOperator :  apply(in,out,n) , dim() =============================
 L4  Operator     CpuMatVecBackend<Policy>   the matvec engine (CSR | matrix-free)
                  CudaMatVecBackend<Policy>
 ===== SEAM:  BasisPolicy : index_of, iter_orbit, coeff_modifier, traits ============
 L5  Basis        Full / FixedSz / Rep /     reduced basis: state -> (index, amplitude)
                  Symmetry / NonAbelianSymmetry
 L6  Sym build    group.cpp, irreps.cpp,     group closure, irrep matrices D^G(g), SAB,
                  symmetry_adapted.cpp        sector packing
 ----------------------------------------------------------------------------------
 L7  Core         core/operator.h            Operator (term SoA), for_each_connected_state
 L8  Compute      term_kernels*, backends/   SpMV kernels; CpuBackend / CudaBackend (BLAS-like)
\end{lstlisting}
\end{center}

The contracts at each seam (Table~\ref{tab:contracts}) are the spine of
the design: every layer talks to its neighbours through exactly one of
these.

\begin{table}[h]
\centering\small
\begin{tabular}{@{}lll@{}}
\toprule
Seam & Contract & Implementors \\
\midrule
L7$\to$L6 & \code{ConnectFn}: \code{connect(s, emit)} $\Rightarrow$ \code{emit(s', $\langle s'|H|s\rangle$)} & \code{Operator::for\_each\_connected\_state} \\
L6$\to$L5 & \code{GroupIrreps} + \code{::SymmetrySector} & \code{decompose\_irreps}, \code{build\_..\_sector} \\
L5$\to$L4 & \code{BasisPolicy} concept (compile-time) & \code{Full/FixedSz/Rep/Symmetry/NonAbelian} \\
L4$\to$L3 & \code{MatVecOperator}: \code{apply}, \code{dim} & \code{CpuMatVecBackend}, \code{CudaMatVecBackend} \\
L3$\to$L2 & method takes \code{const MatVecOperator\&} & \code{lanczos}, \code{full\_diagonalization}, \dots \\
\bottomrule
\end{tabular}
\caption{The five inter-layer contracts.}
\label{tab:contracts}
\end{table}

% =============================================================================
\section{L7 -- Core: the Hamiltonian as a term SoA and a row enumerator}
% =============================================================================

A Hamiltonian is an \code{Operator} (\file{include/ed/core/operator.h}),
which stores terms in a struct-of-arrays (SoA) binned by type
(\file{include/ed/matvec/term\_storage.h}):
\begin{lstlisting}
struct TermStorage {
  std::vector<DiagOneBody>    diag_one_body;     // Sz_i
  std::vector<OffDiagOneBody> offdiag_one_body;  // S+/S-_i
  std::vector<DiagTwoBody>    diag_two_body;     // Sz_i Sz_j
  std::vector<MixedTwoBody>   mixed_two_body;    // Sz_i S+/-_j
  std::vector<OffDiagTwoBody> offdiag_two_body;  // S+/-_i S+/-_j
  std::vector<ThreeBody>      three_body;        // general 3-site
  bool is_real;
};
\end{lstlisting}
The SoA is exposed as a non-owning view \code{TermViewT} (pointers to
the six bins $+$ \code{spin\_l}, \code{is\_real}). Two consumers read it:
\begin{enumerate}[nosep]
  \item the matvec kernels (Section~\ref{sec:kernels}), which apply the
  terms to a basis state;
  \item the \emph{row enumerator}
\begin{lstlisting}
template <class Emit>
void Operator::for_each_connected_state(uint64_t s, Emit&& emit) const;
// invokes emit(s', <s'|H|s>) for every s' connected to s by a term.
\end{lstlisting}
  This is the \code{ConnectFn} contract: a sparse, single-state row of
  $H$ with no $2^N$ vector. The symmetry reduction (L6) uses it to build
  the symmetry-adapted basis, decoupled from how $H$ is stored.
\end{enumerate}

% =============================================================================
\section{L8 -- The rudimentary SpMV kernels}
\label{sec:kernels}
% =============================================================================

The lowest compute layer is a single templated kernel family
(\file{include/ed/matvec/term\_kernels.h} and \code{\_gather.h}),
parameterized on a \code{BasisPolicy} and a \code{Scalar}
(\code{complex<double>} or \code{double}). The structure is: walk a
source column, apply each term type to get a connected state $s'$, and
\emph{emit} the contribution to its destination index. The
\code{emit} lambda is the single funnel where the basis policy enters:

\begin{lstlisting}
auto emit = [&](uint64_t j_idx, uint64_t s_prime, const Scalar& base_contrib) {
  if constexpr (policy_multi_target_v<BasisPolicy>) {
    // Non-abelian: s' belongs to several SAB vectors (multiplicity d_G).
    basis.for_each_target(s_prime, [&](uint64_t k, std::complex<double> w) {
      local_buffer.push_back({k, base_contrib * coerce_coeff<Scalar>(w)});
    });
  } else if constexpr (BasisPolicy::has_coeff_modifier) {
    // Abelian symmetry: one destination, projection phase conj(beta)*group_norm/norm.
    const Scalar mod = basis.template coeff_modifier<Scalar>(basis_state, s_prime, i, j_idx);
    local_buffer.push_back({j_idx, base_contrib * mod});
  } else {
    // Full / FixedSz: identity projection.
    local_buffer.push_back({j_idx, base_contrib});
  }
};
\end{lstlisting}

\noindent
The three branches are selected at compile time by policy traits, so a
policy pays for exactly the work it needs:
\begin{itemize}[nosep]
  \item \textbf{Full / FixedSz}: \code{has\_coeff\_modifier=false},
  \code{multi\_target=false} $\Rightarrow$ the bare push; identical to a
  hand-written SpMV.
  \item \textbf{abelian symmetry} (\code{Rep}/\code{Symmetry}):
  \code{has\_coeff\_modifier=true} $\Rightarrow$ a scalar projection
  phase; one destination per $s'$ (orbits partition the space).
  \item \textbf{non-abelian symmetry}: \code{multi\_target=true}
  $\Rightarrow$ $s'$ maps to \emph{many} destinations (the $d_\Gamma$
  multiplicity copies), enumerated by \code{for\_each\_target}.
\end{itemize}
The detector \code{policy\_multi\_target\_v<P>} defaults to false, so all
pre-existing policies are byte-identical---the multi-target branch is
discarded for them.

There are two kernel forms with identical math: a lock-free row
\textbf{gather} (default; overwrites the output, Hermitian-transpose of
the orbit walk) and a \textbf{scatter} with a radix-sorted accumulation
buffer. The backend chooses; the non-abelian (multi-target) path uses the
scatter form because gather assumes a single destination.

The innermost vector primitives (\code{axpy}, \code{dot}, \code{nrm2})
are a \code{Backend}: \code{CpuBackend} (BLAS) or \code{CudaBackend}
(cuBLAS). This is axis~C in its purest form.

% =============================================================================
\section{L5 -- Basis policies: the reduction as a compile-time policy}
% =============================================================================

A \code{BasisPolicy} maps the full Hilbert space onto a reduced basis.
The contract (a compile-time concept, not an inheritance hierarchy) is:
\code{dim()}, \code{state\_of(i)}, \code{index\_of(s)}, optionally
\code{iter\_orbit} / \code{coeff\_modifier} / \code{for\_each\_target},
plus the boolean traits \code{needs\_orbit\_walk},
\code{has\_coeff\_modifier}, \code{may\_leave\_basis}, \code{multi\_target}.

\begin{table}[h]
\centering\small
\begin{tabular}{@{}lll@{}}
\toprule
Policy & Reduction & Projection \\
\midrule
\code{FullBasisPolicy}        & none ($2^N$)          & identity (bitstring = index) \\
\code{FixedSzBasisPolicy}     & U(1) / fixed $S^z$    & combinadic rank/unrank \\
\code{RepSymmetryBasisPolicy} & abelian spatial       & scalar $\chi_k(g)$, regenerated arithmetically \\
\code{SymmetryBasisPolicy}    & abelian orbit-CSR     & stored orbit coefficients \\
\code{NonAbelianSymmetryBasisPolicy} & non-abelian SAB & stored $D^\Gamma$ amplitudes, multi-target \\
\bottomrule
\end{tabular}
\caption{The five basis policies. All feed the same \code{CpuMatVecBackend}.}
\end{table}

The non-abelian policy
(\file{include/ed/matvec/nonabelian\_symmetry\_basis\_policy.h}) is the
$d_\Gamma\!\ge\!2$ generalization of the abelian one. It views a
\code{::SymmetrySector} of symmetry-adapted basis vectors plus a
multi-target lookup:
\begin{lstlisting}
struct NonAbelianSymmetryBasisPolicy {
  const ::SymmetrySector*      sector;   // SAB vectors (orbit_elements + coeffs)
  const NonAbelianMultiLookup* lookup;   // state -> [basis indices containing it]

  static constexpr bool needs_orbit_walk  = true;
  static constexpr bool has_coeff_modifier = false;  // superseded by for_each_target
  static constexpr bool multi_target      = true;

  int64_t index_of(uint64_t s) const;            // any target, or -1 (in-sector gate)
  template<class Cb> void iter_orbit(uint64_t i, Cb&& cb) const;       // source walk
  template<class Cb> void for_each_target(uint64_t s', Cb&& cb) const; // -> (k, conj(c_k(s')))
};
\end{lstlisting}
Because the SAB is orthonormal ($\text{norm}=1$, $\text{group\_norm}=1$),
the abelian projection formula reduces exactly to $\overline{c_k(s')}$,
so the policy reuses the existing weighting verbatim; only the lookup
becomes multi-valued.

% =============================================================================
\section{L4 -- The matvec engine}
% =============================================================================

\code{CpuMatVecBackend<BasisPolicy, ...terms>}
(\file{include/ed/matvec/matvec\_backend.h}) is the single matvec
engine and a \code{MatVecOperator}. Its \code{apply\_complex(tv, in,
out, n)} dispatches at compile time:
\begin{itemize}[nosep]
  \item \code{needs\_orbit\_walk} (symmetry): always the matrix-free
  orbit-walk kernel. Multi-target $\Rightarrow$ scatter; abelian
  $\Rightarrow$ gather (default) or scatter.
  \item otherwise (Full/FixedSz): \code{csr\_eligible(dim)} chooses an
  assembled CSR SpMV vs.\ a matrix-free SpMV. The choice is made by the
  execution planner (\file{include/ed/planner/}) from a system probe and
  a task cost model, not a hard-coded threshold.
\end{itemize}
\code{CudaMatVecBackend} is the same template specialized for device
memory; it is the axis-C twin.

% =============================================================================
\section{L3 -- Methods: symmetry-blind eigensolvers}
% =============================================================================

Every method consumes only a \code{const MatVecOperator\&}
(\file{include/ed/solvers/lanczos.h}):
\begin{lstlisting}
inline void lanczos(const ed::matvec::MatVecOperator& H, uint64_t N, ...);
inline void full_diagonalization(const ed::matvec::MatVecOperator& H, ...);
inline void krylov_schur(const ed::matvec::MatVecOperator& H, ...);
inline void block_lanczos(const ed::matvec::MatVecOperator& H, ...);
\end{lstlisting}
CPU and GPU Lanczos are the \emph{same} body,
\code{ed::krylov::lanczos\_kernel<Backend>}
(\file{include/ed/krylov/lanczos\_kernel.h}), instantiated with
\code{CpuBackend} or \code{CudaBackend}; the reorthogonalization policy
(None / LocalDGKS3 / PeriodicCGS2 / FullCGS2) is a runtime option. No
method mentions symmetry.

% =============================================================================
\section{L6 -- The symmetry reduction, bottom to top}
% =============================================================================

\subsection{Group closure}
\code{ed::sym::generate\_group(generators)} (\file{src/symmetry/group.cpp})
BFS-closes a set of site permutations into the full group
\code{max\_clique}. A non-abelian guard restricts the abelian-only
\emph{projection} path to a maximal abelian subgroup unless the full
non-abelian machinery is used.

\subsection{Numerical irrep decomposition}
\code{decompose\_irreps(max\_clique, n\_sites)}
(\file{src/symmetry/irreps.cpp}) returns \code{GroupIrreps}: for each
irrep $\Gamma$, the dimension $d_\Gamma$, character $\chi_\Gamma(g)$, and
the $d_\Gamma\!\times\!d_\Gamma$ matrices $D^\Gamma(g)$. The method is a
numerical regular-representation decomposition: form a random Hermitian
element of the group algebra $M=\sum_h c_h R(h)$ (with
$c_{h^{-1}}=\overline{c_h}$), eigendecompose, and read the irreducible
blocks off the eigenspaces. The group product convention
$(g\!\cdot\!h)[i]=\text{perm}_h[\text{perm}_g[i]]$ matches
$U(g)\,|s\rangle = |\,\text{perm}_g(s)\rangle$ so $D^\Gamma$ is matvec-ready.

\subsection{Symmetry-adapted basis (SAB)}
For an irrep $\Gamma$ and partner $p$, the Wigner projector
\[
P^\Gamma_{pp} = \frac{d_\Gamma}{|G|}\sum_{g\in G} \overline{D^\Gamma_{pp}(g)}\,U(g)
\]
applied to an orbit yields $m_\Gamma(O)$ orthonormal partner-$p$ vectors
(\code{build\_sab\_partition0}, \file{src/symmetry/symmetry\_adapted.cpp}),
obtained as the SVD column basis of the projected orbit (absolute
threshold $10^{-8}|G|$ so absent irreps contribute nothing). For a free
orbit, $m_\Gamma=d_\Gamma$, so several SAB vectors share an orbit---this
is the multiplicity that forces the multi-target lookup.

\subsection{Packing into the production sector}
\code{build\_symmetry\_adapted\_sector} packs the partner-$p$ SAB into the
\emph{same} \code{::SymmetrySector} / \code{::SymBasisState} structure
the abelian path uses (\code{orbit\_elements} $=$ states,
\code{orbit\_coefficients} $=$ coeffs, $\text{norm}=1$). An SAB vector
\emph{is} a \code{SymBasisState}; no parallel basis type exists.

% =============================================================================
\section{L2 -- The symmetry-reduced consumers (one engine)}
% =============================================================================

Because each irrep block $H_\Gamma$ is now a \code{MatVecOperator}, the
consumers are thin. \code{ed::solvers::build\_nonabelian\_sector\_matvec}
(\file{src/solvers/cpu/symmetry\_adapted\_solve.cpp}) wraps the sector +
the operator's terms into a \code{CpuMatVecBackend<NonAbelianSymmetry\-BasisPolicy>}
exposed as a \code{MatVecOperator}. The consumers then either feed it to a
method or sample it on the unit columns to materialize $H_\Gamma$ densely:
\begin{lstlisting}
Eigen::MatrixXcd materialize(const MatVecOperator& mv) {  // dense H_Gamma
  const int n = mv.dim();  Eigen::MatrixXcd H(n,n);
  std::vector<Complex> e(n,0), col(n);
  for (int j=0;j<n;++j){ e[j]=1; mv.apply(e.data(),col.data(),n);
    for(int i=0;i<n;++i) H(i,j)=col[i]; e[j]=0; }
  return H;   // one matvec engine, sampled -- not a second materializer
}
\end{lstlisting}
On this base:
\begin{itemize}[nosep]
  \item \code{symmetry\_adapted\_lowest\_eigenvalues} -- per block, dense
  (small) or Lanczos / Krylov--Schur (large) on the engine; the
  \code{BlockMethod} selector makes ``method'' an explicit free choice.
  \item \code{symmetry\_adapted\_full\_spectrum} -- materialize + dense
  eigensolve per block, recombine $\times d_\Gamma$.
  \item \code{symmetry\_adapted\_thermodynamics} -- full spectrum
  $\to$ exact canonical $Z/E/C/S(\beta)$.
  \item \code{symmetry\_adapted\_ground\_state\_dssf} -- per (irrep,
  partner) materialize, eigensolve with eigenvectors, expand to the
  computational basis via the sector, and Lehmann-sum
  $S(\omega)=\sum_n|\langle n|O|0\rangle|^2 L(\omega-(E_n-E_0))$. All
  $d_\Gamma$ partners are summed for completeness.
  \item \code{symmetry\_adapted\_blocks\_packed} -- materialize each block
  column-major and concatenate, for the GPU batched eigensolver.
\end{itemize}

\subsection{GPU}
\code{symmetry\_adapted\_spectrum\_gpu} / \code{\_thermodynamics\_gpu}
(\file{src/solvers/gpu/symmetry\_adapted\_gpu.cu}) take the
\code{::Operator}, materialize the blocks on the host via the CPU engine
(\code{symmetry\_adapted\_blocks\_packed}), upload all blocks in one
transfer, run a multi-stream \code{cusolverDnZheevd} batched eigensolve,
and download the eigenvalues---two host\,$\leftrightarrow$\,device
transfers regardless of block count.

% =============================================================================
\section{End-to-end vertical traces}
% =============================================================================

\paragraph{Ground state, fixed-$S^z$, CPU (production iterative).}
\begin{lstlisting}[language=,basicstyle=\ttfamily\scriptsize]
qed.solve(op, sz=k)                                  python/qed/workflow.py
 -> planner picks matrix-free                        ed/planner/execution_planner.h
 -> lanczos_kernel<CpuBackend>(reorth=FullCGS2)       ed/krylov/lanczos_kernel.h
 -> CpuMatVecBackend<FixedSzBasisPolicy>              ed/matvec/matvec_backend.h
 -> apply_terms over FixedSzOperator term SoA         ed/matvec/term_kernels.h
\end{lstlisting}

\paragraph{Ground state, non-abelian, Krylov--Schur, CPU.}
\begin{lstlisting}[language=,basicstyle=\ttfamily\scriptsize]
qed._core.symmetry_adapted_lowest(op, gens, method)  qed_bindings.cpp
 -> symmetry_adapted_lowest_eigenvalues               ed/solvers/symmetry_adapted_solve.*
      decompose_irreps(generate_group(gens))          ed/symmetry/irreps.cpp, group.cpp
      build_symmetry_adapted_sector                   ed/symmetry/symmetry_adapted.cpp
      CpuMatVecBackend<NonAbelianSymmetryBasisPolicy>  ed/matvec/matvec_backend.h  (= engine)
      krylov_schur(mv, n_G, ...)                       ed/solvers/lanczos.h
       -> apply_terms (scatter, multi-target emit)     ed/matvec/term_kernels.h
\end{lstlisting}

\paragraph{Finite-$T$, non-abelian, GPU.}
\begin{lstlisting}[language=,basicstyle=\ttfamily\scriptsize]
qed._core.symmetry_adapted_thermodynamics(op,gens,T,use_gpu=True)
 -> symmetry_adapted_thermodynamics_gpu(op,...)        symmetry_adapted_gpu.cu
      symmetry_adapted_blocks_packed(op,...)           (CPU engine materializes blocks)
      upload -> multi-stream cusolverDnZheevd -> down  (GPU batched eigensolve)
      canonical_thermo_from_eigs                       ed/symmetry/symmetry_adapted.cpp
\end{lstlisting}

% =============================================================================
\section{Coverage matrix (verified)}
% =============================================================================

Every cell below is validated against brute-force full diagonalization
(or, for GPU, against the CPU engine) to $\sim 10^{-15}$ on the $N=6$
Heisenberg ring (\file{tests/unit/test\_symmetry\_adapted.cpp} and the
Python end-to-end harness).

\begin{table}[h]
\centering\small
\begin{tabular}{@{}lcccc@{}}
\toprule
Reduction & dense & Lanczos / KS & finite-$T$ & DSSF \\
\midrule
Full / Sz            & CPU            & CPU$\cdot$GPU$\cdot$MPI & CPU$\cdot$GPU & CPU \\
abelian spatial      & CPU            & CPU$\cdot$GPU$\cdot$MPI & CPU & CPU \\
non-abelian          & CPU$\cdot$GPU  & CPU            & CPU$\cdot$GPU & CPU \\
$S^z\times$ symmetry & CPU$\cdot$GPU  & CPU            & CPU$\cdot$GPU & CPU \\
\bottomrule
\end{tabular}
\caption{Reduction $\times$ method $\times$ deployment. All cells run on the
single \code{CpuMatVecBackend} / \code{CudaMatVecBackend} engine.}
\end{table}

% =============================================================================
\section{Summary}
% =============================================================================

The architecture is one engine and three orthogonal axes. A Hamiltonian
is a term SoA with a row enumerator (L7). A reduction is a basis policy
(L5) built by the group/irrep/SAB machinery (L6). The engine
\code{CpuMatVecBackend<Policy>} (L4) turns (terms $+$ policy) into a
\code{MatVecOperator}. A method (L3) consumes that operator, and a
backend (L8) supplies the vector primitives---CPU, GPU, or MPI. Non-abelian
symmetry, historically a parallel subsystem, is now just another basis
policy on the same engine: the only thing it needed beyond the abelian
case was a multi-valued lookup ($s' \mapsto$ several SAB vectors), added
as a single compile-time-gated branch in the SpMV \code{emit}.

\end{document}
